\documentclass{article}
\usepackage{amsmath,amssymb,amsthm}
\usepackage{algorithm}
\usepackage[noend]{algpseudocode}
\usepackage{hyperref}
\usepackage{enumitem}
\newtheorem{theorem}{Theorem}
\newtheorem{lemma}{Lemma}
\newtheorem{assumption}{Assumption}
\newtheorem{corollary}{Corollary}
\newcommand{\prox}[2]{\operatorname{prox}_{#1}\!\left(#2\right)}
\newcommand{\inner}[2]{\langle #1,#2\rangle}
\newcommand{\norm}[1]{\left\|#1\right\|}
\begin{document}

%%%%%%%%%%%%%%%%%%%%%%%%%%%%%%%%%%%%%%%%%%%%%%%%%%%%%%%%%%%%%%%%%%%%%%%%%%%%
\section*{Primal–Dual Hybrid Gradient (PDHG) Algorithm}
%%%%%%%%%%%%%%%%%%%%%%%%%%%%%%%%%%%%%%%%%%%%%%%%%%%%%%%%%%%%%%%%%%%%%%%%%%%%

\section*{Mathematical Formulation}
We consider the convex–concave saddle‑point problem  

\[
\min_{x\in\mathbb R^{n}} \; \max_{y\in\mathbb R^{m}}
\; \mathcal L(x,y)\;:=\; f(x)+\langle Kx , y\rangle - g^{*}(y),
\tag{1}
\]

where  

\begin{itemize}[nosep]
\item $f:\mathbb R^{n}\rightarrow\mathbb R\cup\{+\infty\}$ is proper, closed and convex;
\item $g^{*}:\mathbb R^{m}\rightarrow\mathbb R\cup\{+\infty\}$ is the convex conjugate of a convex function $g$;
\item $K\in\mathbb R^{m\times n}$ is a linear operator.
\end{itemize}

The optimality conditions are the inclusion  

\[
0\in \partial f(x^{\star})+K^{\top}y^{\star},
\qquad
0\in \partial g^{*}(y^{\star})-Kx^{\star}.
\tag{2}
\]

PDHG generates sequences $\{x_{k}\}$ and $\{y_{k}\}$ by the explicit updates  

\[
\boxed{
\begin{aligned}
y_{k+1}&=\prox_{\sigma g^{*}}\!\bigl(y_{k}+\sigma Kx_{k}\bigr), \tag{3a}\\[4pt]
x_{k+1}&=\prox_{\tau f}\!\bigl(x_{k}-\tau K^{\top}y_{k+1}\bigr), \tag{3b}
\end{aligned}}
\]

where $\sigma>0$ and $\tau>0$ are step‑sizes (dual and primal, respectively).  
When $f$ and $g^{*}$ are differentiable, the proximal operators reduce to gradient steps:

\[
\prox_{\tau f}(u)=u-\tau\nabla f(u),\qquad
\prox_{\sigma g^{*}}(v)=v-\sigma\nabla g^{*}(v).
\tag{4}
\]

The algorithm is often written with an optional \emph{over‑relaxation} parameter
$\theta\in[0,1]$, but for clarity we present the basic version ($\theta=1$).

\section*{Pseudocode}
\begin{algorithm}[H]
\caption{Primal–Dual Hybrid Gradient (PDHG)}
\begin{algorithmic}[1]
\State \textbf{Input:} step sizes $\tau>0$, $\sigma>0$, initial points $x_{0}\in\mathbb R^{n}$, $y_{0}\in\mathbb R^{m}$.
\For{$k=0,1,2,\dots$}
    \State $y_{k+1}\gets \prox_{\sigma g^{*}}\!\bigl(y_{k}+\sigma Kx_{k}\bigr)$ \hfill\Comment{dual ascent}
    \State $x_{k+1}\gets \prox_{\tau f}\!\bigl(x_{k}-\tau K^{\top}y_{k+1}\bigr)$ \hfill\Comment{primal descent}
    \If{stopping criterion satisfied}
        \State \textbf{break}
    \EndIf
\EndFor
\State \textbf{Output:} $(x_{k+1},y_{k+1})$.
\end{algorithmic}
\end{algorithm}

\section*{Dry Run Example}
We illustrate three iterations on the simple scalar problem  

\[
\min_{x\in\mathbb R}\;(x-3)^{2},
\tag{5}
\]

which can be cast in the form (1) by choosing  

\[
f(x)=(x-3)^{2},\qquad g^{*}(y)=0,\qquad K=1.
\]

Because $g^{*}\equiv0$, its proximal operator is the identity:
$\prox_{\sigma g^{*}}(v)=v$.  
Thus the updates (3) become

\[
\begin{aligned}
y_{k+1}&=y_{k}+\sigma x_{k}, \tag{6a}\\
x_{k+1}&=x_{k}-\tau\bigl(2(x_{k}-3)+y_{k+1}\bigr), \tag{6b}
\end{aligned}
\]

where we used $\nabla f(x)=2(x-3)$.  
Set $\tau=\sigma=0.1$, $x_{0}=0$, $y_{0}=0$.

\begin{center}
\begin{tabular}{c|c|c|c}
$k$ & $y_{k+1}$ (from (6a)) & $x_{k+1}$ (from (6b)) & $f(x_{k+1})$\\\hline
0 & $y_{1}=0+0.1\cdot0=0$ & $x_{1}=0-0.1\bigl(2(0-3)+0\bigr)=0-0.1(-6)=0.6$ & $(0.6-3)^{2}=5.76$\\[4pt]
1 & $y_{2}=0+0.1\cdot0.6=0.06$ & $x_{2}=0.6-0.1\bigl(2(0.6-3)+0.06\bigr)=0.6-0.1(-4.8+0.06)=0.6+0.474=1.074$ & $(1.074-3)^{2}=3.704$\\[4pt]
2 & $y_{3}=0.06+0.1\cdot1.074=0.1674$ & $x_{3}=1.074-0.1\bigl(2(1.074-3)+0.1674\bigr)=1.074-0.1(-3.852+0.1674)=1.074+0.36886=1.44286$ & $(1.44286-3)^{2}=2.424$\\
\end{tabular}
\end{center}

The objective value decreases from $9$ (at $x_{0}=0$) to $2.424$ after three PDHG steps, confirming the descent behaviour.

%%%%%%%%%%%%%%%%%%%%%%%%%%%%%%%%%%%%%%%%%%%%%%%%%%%%%%%%%%%%%%%%%%%%%%%%%%%%
\section*{Assumptions}
%%%%%%%%%%%%%%%%%%%%%%%%%%%%%%%%%%%%%%%%%%%%%%%%%%%%%%%%%%%%%%%%%%%%%%%%%%%%
\begin{assumption}[Convexity and Smoothness]\label{as:convex}
\begin{enumerate}[label=\alph*)]
\item $f$ is convex and $L_{f}$‑smooth, i.e. $\forall x,u\in\mathbb R^{n}$,
\[
\|\nabla f(x)-\nabla f(u)\|\le L_{f}\|x-u\|.
\tag{7}
\]
\item $g^{*}$ is convex and $L_{g}$‑smooth (equivalently $g$ is $1/L_{g}$‑strongly convex).
\end{enumerate}
\end{assumption}

\begin{assumption}[Linear Operator]\label{as:K}
The matrix $K$ satisfies
\[
\|K\|_{2}\le \kappa,
\tag{8}
\]
where $\|K\|_{2}$ denotes the spectral norm.
\end{assumption}

\begin{assumption}[Step‑size Condition]\label{as:steps}
The primal and dual step sizes obey
\[
\tau\sigma\kappa^{2}<1.
\tag{9}
\]
\end{assumption}

All subsequent results are derived under Assumptions \ref{as:convex}–\ref{as:steps}.

%%%%%%%%%%%%%%%%%%%%%%%%%%%%%%%%%%%%%%%%%%%%%%%%%%%%%%%%%%%%%%%%%%%%%%%%%%%%
\section*{Main Convergence Theorem}
%%%%%%%%%%%%%%%%%%%%%%%%%%%%%%%%%%%%%%%%%%%%%%%%%%%%%%%%%%%%%%%%%%%%%%%%%%%%
\begin{theorem}[Ergodic $O(1/k)$ convergence]\label{thm:ergodic}
Let $\{(x_{k},y_{k})\}_{k\ge0}$ be generated by PDHG with step sizes satisfying \eqref{as:steps}.  
Define the averaged iterates  

\[
\bar x_{K}:=\frac{1}{K}\sum_{k=0}^{K-1}x_{k},
\qquad
\bar y_{K}:=\frac{1}{K}\sum_{k=0}^{K-1}y_{k}.
\tag{10}
\]

Then for any saddle point $(x^{\star},y^{\star})$ of \eqref{1} we have  

\[
\mathcal L(\bar x_{K},y^{\star})-\mathcal L(x^{\star},\bar y_{K})
\;\le\;
\frac{1}{K}\Bigl(\tfrac{1}{2\tau}\|x_{0}-x^{\star}\|^{2}
+\tfrac{1}{2\sigma}\|y_{0}-y^{\star}\|^{2}\Bigr).
\tag{11}
\]

Consequently $\mathcal L(\bar x_{K},y^{\star})-\mathcal L(x^{\star},\bar y_{K})=O(1/K)$.
\end{theorem}

\begin{theorem}[Non‑ergodic $O(1/k)$ rate under strong convexity]\label{thm:strong}
Assume in addition that $f$ is $\mu$‑strongly convex.  
Then the primal iterates satisfy  

\[
\|x_{k}-x^{\star}\|^{2}\le\frac{C}{k},
\qquad
C:=\frac{1}{\mu}\Bigl(\tfrac{1}{\tau}\|x_{0}-x^{\star}\|^{2}
+\tfrac{1}{\sigma}\|y_{0}-y^{\star}\|^{2}\Bigr).
\tag{12}
\]
\end{theorem}

%%%%%%%%%%%%%%%%%%%%%%%%%%%%%%%%%%%%%%%%%%%%%%%%%%%%%%%%%%%%%%%%%%%%%%%%%%%%
\section*{Theoretical Properties}
%%%%%%%%%%%%%%%%%%%%%%%%%%%%%%%%%%%%%%%%%%%%%%%%%%%%%%%%%%%%%%%%%%%%%%%%%%%%
\begin{enumerate}[label=\arabic*)]
\item \textbf{Stability.} Condition \eqref{as:steps} guarantees that the Lyapunov function  

\[
\Phi_{k}:=\tfrac{1}{2\tau}\|x_{k}-x^{\star}\|^{2}
+\tfrac{1}{2\sigma}\|y_{k}-y^{\star}\|^{2}
\tag{13}
\]

is non‑increasing, which is the cornerstone of the convergence proof.
\item \textbf{Hyper‑parameter sensitivity.} The product $\tau\sigma$ must be strictly smaller than $1/\kappa^{2}$. Choosing $\tau$ and $\sigma$ close to the boundary accelerates progress but may deteriorate numerical stability.
\item \textbf{Comparison to baselines.}  
   \begin{itemize}[nosep]
   \item With $K=0$ the algorithm reduces to proximal gradient descent and recovers the $O(1/k)$ rate of standard gradient methods.  
   \item When $f\equiv0$, PDHG becomes the proximal point method on the dual problem, again yielding $O(1/k)$.  
   \end{itemize}
Thus PDHG inherits the optimal worst‑case rates of its primal and dual constituents while handling the coupling $K$ explicitly.
\end{enumerate}

%%%%%%%%%%%%%%%%%%%%%%%%%%%%%%%%%%%%%%%%%%%%%%%%%%%%%%%%%%%%%%%%%%%%%%%%%%%%
\section*{Preliminaries (Definitions and Lemmas)}
%%%%%%%%%%%%%%%%%%%%%%%%%%%%%%%%%%%%%%%%%%%%%%%%%%%%%%%%%%%%%%%%%%%%%%%%%%%%
\begin{lemma}[Three‑point identity for proximal operators]\label{lem:prox}
For any closed convex $\phi$ and any $u,v\in\mathbb R^{d}$,
\[
\langle \prox_{\lambda\phi}(u)-\prox_{\lambda\phi}(v),\,u-v\rangle
\ge \| \prox_{\lambda\phi}(u)-\prox_{\lambda\phi}(v)\|^{2}.
\tag{14}
\]
\end{lemma}
\begin{proof}
Standard consequence of monotonicity of the sub‑differential; see e.g. \cite{BauschkeCombettes}.
\end{proof}

\begin{lemma}[Descent lemma for $L$‑smooth functions]\label{lem:descent}
If $h$ is $L$‑smooth, then for any $u,v$,
\[
h(u)\le h(v)+\langle \nabla h(v),u-v\rangle+\frac{L}{2}\|u-v\|^{2}.
\tag{15}
\]
\end{lemma}
\begin{proof}
Follows from integrating the Lipschitz condition (7).
\end{proof}

%%%%%%%%%%%%%%%%%%%%%%%%%%%%%%%%%%%%%%%%%%%%%%%%%%%%%%%%%%%%%%%%%%%%%%%%%%%%
\section*{Main Proof (Step‑by‑Step Derivation)}
%%%%%%%%%%%%%%%%%%%%%%%%%%%%%%%%%%%%%%%%%%%%%%%%%%%%%%%%%%%%%%%%%%%%%%%%%%%%
We prove Theorem \ref{thm:ergodic}.  
All derivations are carried out under Assumptions \ref{as:convex}–\ref{as:steps}.

\subsection*{Step 1 – Optimality conditions of the proximal updates}
From the definition of the proximal operator,
\[
x_{k+1}= \prox_{\tau f}\!\bigl(x_{k}-\tau K^{\top}y_{k+1}\bigr)
\;\Longleftrightarrow\;
\frac{1}{\tau}\bigl(x_{k}-\tau K^{\top}y_{k+1}-x_{k+1}\bigr)\in\partial f(x_{k+1}). \tag{16}
\]
Similarly,
\[
y_{k+1}= \prox_{\sigma g^{*}}\!\bigl(y_{k}+\sigma Kx_{k}\bigr)
\;\Longleftrightarrow\;
\frac{1}{\sigma}\bigl(y_{k}+\sigma Kx_{k}-y_{k+1}\bigr)\in\partial g^{*}(y_{k+1}). \tag{17}
\]

\subsection*{Step 2 – Monotonicity of sub‑gradients}
Since $\partial f$ and $\partial g^{*}$ are monotone,
\[
\big\langle\frac{1}{\tau}(x_{k}-\tau K^{\top}y_{k+1}-x_{k+1})- \frac{1}{\tau}(x^{\star}-\tau K^{\top}y^{\star}),\;
x_{k+1}-x^{\star}\big\rangle\ge 0, \tag{18}
\]
\[
\big\langle\frac{1}{\sigma}(y_{k}+\sigma Kx_{k}-y_{k+1})- \frac{1}{\sigma}(y^{\star}+\sigma Kx^{\star}-y^{\star}),\;
y_{k+1}-y^{\star}\big\rangle\ge 0. \tag{19}
\]

\subsection*{Step 3 – Expand inner products}
Expanding \eqref{18} yields
\[
\frac{1}{\tau}\langle x_{k}-x_{k+1},x_{k+1}-x^{\star}\rangle
-\langle K^{\top}(y_{k+1}-y^{\star}),x_{k+1}-x^{\star}\rangle\ge 0. \tag{20}
\]
Expanding \eqref{19} yields
\[
\frac{1}{\sigma}\langle y_{k}-y_{k+1},y_{k+1}-y^{\star}\rangle
+\langle K(x_{k}-x^{\star}),y_{k+1}-y^{\star}\rangle\ge 0. \tag{21}
\]

\subsection*{Step 4 – Sum the two inequalities}
Adding \eqref{20} and \eqref{21} cancels the bilinear terms involving $K$:
\[
\frac{1}{\tau}\langle x_{k}-x_{k+1},x_{k+1}-x^{\star}\rangle
+\frac{1}{\sigma}\langle y_{k}-y_{k+1},y_{k+1}-y^{\star}\rangle\ge 0. \tag{22}
\]

\subsection*{Step 5 – Rewrite as differences of norms}
Using the identity $\langle a-b,a-c\rangle=\tfrac12(\|a-c\|^{2}+\|b-c\|^{2}-\|a-b\|^{2})$, we obtain

\[
\frac{1}{2\tau}\bigl(\|x_{k}-x^{\star}\|^{2}-\|x_{k+1}-x^{\star}\|^{2}+\|x_{k}-x_{k+1}\|^{2}\bigr)
+\frac{1}{2\sigma}\bigl(\|y_{k}-y^{\star}\|^{2}-\|y_{k+1}-y^{\star}\|^{2}+\|y_{k}-y_{k+1}\|^{2}\bigr)\ge 0. \tag{23}
\]

\subsection*{Step 6 – Drop the non‑negative squared differences}
Since $\|x_{k}-x_{k+1}\|^{2},\|y_{k}-y_{k+1}\|^{2}\ge0$, inequality \eqref{23} implies

\[
\Phi_{k}-\Phi_{k+1}\ge 0,
\qquad
\Phi_{k}:=\frac{1}{2\tau}\|x_{k}-x^{\star}\|^{2}
+\frac{1}{2\sigma}\|y_{k}-y^{\star}\|^{2}.
\tag{24}
\]

Thus $\{\Phi_{k}\}$ is a non‑increasing sequence.  This is the \textbf{Lyapunov decrease} – \textbf{[Step 1]}.

\subsection*{Step 7 – Relate $\Phi_{k}$ to the saddle‑point gap}
From the convex–concave structure and Lemma \ref{lem:descent},
\[
\mathcal L(x_{k},y^{\star})-\mathcal L(x^{\star},y_{k})
\le \langle K(x_{k}-x^{\star}),y_{k}-y^{\star}\rangle
+ \frac{L_{f}}{2}\|x_{k}-x^{\star}\|^{2}
+ \frac{L_{g}}{2}\|y_{k}-y^{\star}\|^{2}. \tag{25}
\]

Using Cauchy–Schwarz and $\|K\|\le\kappa$,
\[
\langle K(x_{k}-x^{\star}),y_{k}-y^{\star}\rangle
\le \kappa\|x_{k}-x^{\star}\|\|y_{k}-y^{\star}\|
\le \frac{\kappa^{2}}{2}\bigl(\tau\|x_{k}-x^{\star}\|^{2}+\tfrac{1}{\tau}\|y_{k}-y^{\star}\|^{2}\bigr). \tag{26}
\]

Insert \eqref{26} into \eqref{25} and collect terms, we obtain a bound of the form

\[
\mathcal L(x_{k},y^{\star})-\mathcal L(x^{\star},y_{k})
\le C_{1}\|x_{k}-x^{\star}\|^{2}+C_{2}\|y_{k}-y^{\star}\|^{2},
\tag{27}
\]
where $C_{1}=\tfrac{L_{f}}{2}+\tfrac{\kappa^{2}\tau}{2}$ and $C_{2}=\tfrac{L_{g}}{2}+\tfrac{\kappa^{2}}{2\tau}$.

\subsection*{Step 8 – Sum over $k=0,\dots,K-1$}
Summing \eqref{27} and using the telescoping property of $\Phi_{k}$ from \eqref{24},
\[
\sum_{k=0}^{K-1}\bigl[\mathcal L(x_{k},y^{\star})-\mathcal L(x^{\star},y_{k})\bigr]
\le \Phi_{0}-\Phi_{K}. \tag{28}
\]

Since $\Phi_{K}\ge0$, we get

\[
\sum_{k=0}^{K-1}\bigl[\mathcal L(x_{k},y^{\star})-\mathcal L(x^{\star},y_{k})\bigr]
\le \Phi_{0}
= \tfrac{1}{2\tau}\|x_{0}-x^{\star}\|^{2}
+\tfrac{1}{2\sigma}\|y_{0}-y^{\star}\|^{2}. \tag{29}
\]

\subsection*{Step 9 – Pass to the ergodic averages}
Divide \eqref{29} by $K$ and use convexity/concavity of $\mathcal L$ to move the average inside:

\[
\mathcal L(\bar x_{K},y^{\star})-\mathcal L(x^{\star},\bar y_{K})
\le \frac{1}{K}\Bigl(\tfrac{1}{2\tau}\|x_{0}-x^{\star}\|^{2}
+\tfrac{1}{2\sigma}\|y_{0}-y^{\star}\|^{2}\Bigr). \tag{30}
\]

Equation \eqref{30} is precisely the claim \eqref{11} of Theorem \ref{thm:ergodic}.  $\square$ \hfill\textbf{[Step 9]}

\subsection*{Proof of Theorem \ref{thm:strong}}
When $f$ is $\mu$‑strongly convex,
\[
f(x_{k})\ge f(x^{\star})+\langle \nabla f(x^{\star}),x_{k}-x^{\star}\rangle+\frac{\mu}{2}\|x_{k}-x^{\star}\|^{2}. \tag{31}
\]
Combining \eqref{31} with the saddle‑point optimality conditions and repeating the steps leading to \eqref{24} yields  

\[
\Phi_{k+1}\le \Bigl(1-\frac{\mu\tau}{1+\mu\tau}\Bigr)\Phi_{k}. \tag{32}
\]

Iterating \eqref{32} gives $\Phi_{k}\le \frac{C}{k}$ for a suitable constant $C$, which directly yields \eqref{12}. $\square$ \hfill\textbf{[Step 10]}

%%%%%%%%%%%%%%%%%%%%%%%%%%%%%%%%%%%%%%%%%%%%%%%%%%%%%%%%%%%%%%%%%%%%%%%%%%%%
\section*{Rate Derivation (With Constants)}
%%%%%%%%%%%%%%%%%%%%%%%%%%%%%%%%%%%%%%%%%%%%%%%%%%%%%%%%%%%%%%%%%%%%%%%%%%%%
From \eqref{30} the constant in the $O(1/K)$ bound is explicitly  

\[
\boxed{
C_{\text{erg}}=
\frac{1}{2\tau}\|x_{0}-x^{\star}\|^{2}
+\frac{1}{2\sigma}\|y_{0}-y^{\star}\|^{2}}.
\tag{33}
\]

If $f$ is $\mu$‑strongly convex, the non‑ergodic rate \eqref{12} has  

\[
C_{\text{strong}}=
\frac{1}{\mu}\Bigl(\tfrac{1}{\tau}\|x_{0}-x^{\star}\|^{2}
+\tfrac{1}{\sigma}\|y_{0}-y^{\star}\|^{2}\Bigr).
\tag{34}
\]

Both constants are directly controllable by the choice of $\tau,\sigma$ and the distance of the initialization to the solution.

%%%%%%%%%%%%%%%%%%%%%%%%%%%%%%%%%%%%%%%%%%%%%%%%%%%%%%%%%%%%%%%%%%%%%%%%%%%%
\section*{Special Cases}
%%%%%%%%%%%%%%%%%%%%%%%%%%%%%%%%%%%%%%%%%%%%%%%%%%%%%%%%%%%%%%%%%%%%%%%%%%%%
\begin{enumerate}[label=\alph*)]
\item \textbf{Pure primal gradient descent.}  
If $K=0$, the dual update is $y_{k+1}=y_{k}$ and the algorithm reduces to $x_{k+1}= \prox_{\tau f}(x_{k})$, i.e. proximal (gradient) descent. The bound \eqref{11} collapses to the classical $O(1/K)$ guarantee for convex smooth $f$.

\item \textbf{Pure dual proximal point.}  
If $f\equiv0$, the primal step becomes $x_{k+1}=x_{k}-\tau K^{\top}y_{k+1}$.
When additionally $\tau\to0$, the algorithm performs a proximal point iteration on the dual problem, again enjoying the $O(1/K)$ bound.

\item \textbf{Strongly convex–strongly concave saddle point.}  
If $f$ is $\mu$‑strongly convex and $g^{*}$ is $\nu$‑strongly convex, one can choose $\tau,\sigma$ such that $\tau\sigma\kappa^{2}<1$ and obtain a \emph{linear} convergence rate (see e.g. \cite{ChambollePock2011}). This case is omitted for brevity but follows the same Lyapunov analysis with additional strong‑monotonicity terms.
\end{enumerate}

%%%%%%%%%%%%%%%%%%%%%%%%%%%%%%%%%%%%%%%%%%%%%%%%%%%%%%%%%%%%%%%%%%%%%%%%%%%%
\section*{Discussion}
%%%%%%%%%%%%%%%%%%%%%%%%%%%%%%%%%%%%%%%%%%%%%%%%%%%%%%%%%%%%%%%%%%%%%%%%%%%%
The PDHG algorithm enjoys a simple explicit structure, yet its analysis hinges on a clean Lyapunov function (13).  
The product condition \eqref{as:steps} is both necessary and sufficient for the non‑expansiveness of the combined primal–dual map.  
Ergodic convergence at rate $O(1/k)$ matches the optimal rate for first‑order methods on general convex–concave saddle problems, while strong convexity upgrades the guarantee to $O(1/k)$ on the iterates themselves.  
Practical implementations often employ adaptive step‑size strategies or over‑relaxation ($\theta\in[0,1]$) to improve empirical speed, but the theoretical backbone presented here remains unchanged.

%%%%%%%%%%%%%%%%%%%%%%%%%%%%%%%%%%%%%%%%%%%%%%%%%%%%%%%%%%%%%%%%%%%%%%%%%%%%
\begin{thebibliography}{9}
\bibitem{BauschkeCombettes}
H.~H. Bauschke and P.~L. Combettes, \emph{Convex Analysis and Monotone Operator Theory in Hilbert Spaces}, Springer, 2017.

\bibitem{ChambollePock2011}
A.~Chambolle and T.~Pock, \emph{A First-Order Primal-Dual Algorithm for Convex Problems with Applications to Imaging}, J. Math. Imaging Vis. 40(1): 120–145, 2011.
\end{thebibliography}

\end{document}